%%=============================================================================
%% Conclusie
%%=============================================================================

\chapter{Conclusie}%
\label{ch:conclusie}

% TODO: Trek een duidelijke conclusie, in de vorm van een antwoord op de
% onderzoeksvra(a)g(en). Wat was jouw bijdrage aan het onderzoeksdomein en
% hoe biedt dit meerwaarde aan het vakgebied/doelgroep? 
% Reflecteer kritisch over het resultaat. In Engelse teksten wordt deze sectie
% ``Discussion'' genoemd. Had je deze uitkomst verwacht? Zijn er zaken die nog
% niet duidelijk zijn?
% Heeft het onderzoek geleid tot nieuwe vragen die uitnodigen tot verder 
%onderzoek?

% TODO: Schrijf de conclusie nadat alle resultaten en de vergelijkende analyse
% (ch:vergelijkende-analyse) zijn afgerond. Behandel de volgende punten:
%
%   1. Antwoord op de centrale onderzoeksvraag:
%      Wat is het effect van vroeg, laat of gecombineerd testen op de
%      gebruiksvriendelijkheid van een React-applicatie?
%
%   2. Bespreking van de vier hypothesen (H0,1 t/m H0,4):
%      Welke werden verworpen? Welke niet? Wat betekent dit?
%
%   3. Praktische aanbevelingen voor het werkveld:
%      Wanneer is usability-testing het meest effectief?
%      Wat levert telemetry toe bovenop klassieke tests?
%
%   4. Beperkingen van het onderzoek:
%      - Kleine steekproefgroottes (n=13-14 per groep)
%      - Drie verschillende formulierthema's (niet volledig equivalent)
%      - Zelfselectie van deelnemers
%
%   5. Suggesties voor vervolgonderzoek:
%      - Grotere steekproef
%      - Andere types applicaties (niet enkel formulieren)
%      - Longitudinaal onderzoek over meerdere sprints

